% ============================================================
%  MyHR — Graduation Documentation
%  Ain Shams University · Faculty of Computer & Information Sciences
%  Compile with: pdflatex main.tex  (run twice for TOC/LOF/LOT)
% ============================================================
\documentclass[12pt,a4paper,oneside]{report}

% ── Encoding & font ──────────────────────────────────────────
\usepackage[utf8]{inputenc}
\usepackage[T1]{fontenc}
\usepackage{mathptmx}          % Times New Roman (text + math)

% ── Page geometry ────────────────────────────────────────────
\usepackage[
  top=2.5cm, bottom=2.5cm,
  left=3.0cm, right=2.5cm
]{geometry}

% ── Line spacing ─────────────────────────────────────────────
\usepackage{setspace}
\onehalfspacing

% ── Paragraph style ──────────────────────────────────────────
\usepackage{parskip}
\setlength{\parindent}{0pt}
\setlength{\parskip}{6pt plus 2pt minus 1pt}

% ── Colours ──────────────────────────────────────────────────
\usepackage[dvipsnames,svgnames,table]{xcolor}
\definecolor{ainshams}{RGB}{143,0,0}    % Ain Shams dark red
\definecolor{linkblue}{RGB}{0,70,140}

% ── Headers / footers ────────────────────────────────────────
\usepackage{fancyhdr}
\pagestyle{fancy}
\fancyhf{}
\fancyhead[L]{\small\leftmark}
\fancyhead[R]{\small MyHR}
\fancyfoot[C]{\small\thepage}
\renewcommand{\headrulewidth}{0.4pt}
\renewcommand{\footrulewidth}{0pt}

% ── Chapter / section title formatting ───────────────────────
\usepackage{titlesec}
\titleformat{\chapter}[display]
  {\normalfont\huge\bfseries\color{ainshams}}
  {\chaptertitlename\ \thechapter}{18pt}{\Huge}
\titlespacing*{\chapter}{0pt}{-10pt}{20pt}

\titleformat{\section}
  {\normalfont\Large\bfseries\color{ainshams}}
  {\thesection}{1em}{}
\titleformat{\subsection}
  {\normalfont\large\bfseries}{\thesubsection}{1em}{}
\titleformat{\subsubsection}
  {\normalfont\normalsize\bfseries}{\thesubsubsection}{1em}{}

% ── TOC / LOF / LOT styling ──────────────────────────────────
\usepackage{tocloft}
\renewcommand{\cftchapfont}{\bfseries}
\renewcommand{\cftchappagefont}{\bfseries}
\setlength{\cftbeforechapskip}{6pt}

% ── Graphics ─────────────────────────────────────────────────
\usepackage{graphicx}
\usepackage{float}
\usepackage{caption}
\captionsetup{font=small, labelfont=bf, labelsep=period}

% ── Figure / table numbering by chapter ──────────────────────
\usepackage{chngcntr}
\counterwithin{figure}{chapter}
\counterwithin{table}{chapter}

% ── Tables ───────────────────────────────────────────────────
\usepackage{booktabs}
\usepackage{longtable}
\usepackage{tabularx}
\usepackage{multirow}
\usepackage{array}
\newcolumntype{L}[1]{>{\raggedright\arraybackslash}p{#1}}
\newcolumntype{C}[1]{>{\centering\arraybackslash}p{#1}}

% ── Lists ────────────────────────────────────────────────────
\usepackage{enumitem}
\setlist{nosep, topsep=2pt}

% ── Code listings ────────────────────────────────────────────
\usepackage{listings}
\lstset{
  basicstyle=\ttfamily\small,
  breaklines=true,
  frame=single,
  backgroundcolor=\color{gray!8},
  rulecolor=\color{gray!40},
  aboveskip=8pt, belowskip=8pt,
  xleftmargin=10pt, xrightmargin=10pt
}

% ── Diagram placeholder helper ────────────────────────────────
\usepackage{tcolorbox}
\tcbuselibrary{skins}
\newcommand{\diagramplaceholder}[1]{%
  \begin{tcolorbox}[
    colback=gray!5, colframe=gray!50,
    boxrule=0.4pt, arc=2pt,
    left=10pt, right=10pt, top=8pt, bottom=8pt,
    fontupper=\small\itshape
  ]
  \centering #1\\[4pt]
  \textbf{[Replace with exported diagram image: \texttt{\textbackslash includegraphics\{...\}}]}
  \end{tcolorbox}%
}

% ── Hyperlinks (load last) ────────────────────────────────────
\usepackage[
  colorlinks=true,
  linkcolor=linkblue,
  citecolor=linkblue,
  urlcolor=linkblue,
  bookmarks=true,
  bookmarksnumbered=true,
  pdftitle={MyHR -- An Adaptive AI-Powered Hiring \& Interview Platform},
  pdfauthor={Ain Shams University -- Computer Science},
  pdfsubject={Graduation Project Documentation}
]{hyperref}
\usepackage{bookmark}

% ── Chapter intro-page command ────────────────────────────────
% Usage: \chapterfrontpage{Chapter One}{Introduction}{outline items}
\newcommand{\chapterfrontpage}[3]{%
  \clearpage
  \thispagestyle{empty}
  \vspace*{4cm}
  \begin{center}
    \textcolor{ainshams}{\rule{0.6\linewidth}{1.5pt}}\\[0.6cm]
    {\LARGE\textsc{#1}}\\[0.8cm]
    {\Huge\bfseries #2}\\[0.6cm]
    \textcolor{ainshams}{\rule{0.6\linewidth}{1.5pt}}
  \end{center}
  \vspace{1.8cm}
  \begin{center}
    \begin{minipage}{0.72\linewidth}
      \textbf{\large Chapter Outline}\\[6pt]
      \begin{itemize}[leftmargin=1.4em]
        #3
      \end{itemize}
    \end{minipage}
  \end{center}
  \clearpage
}

% ── Misc ─────────────────────────────────────────────────────
\usepackage{microtype}
\usepackage{amsmath}
\usepackage{csquotes}

% ─────────────────────────────────────────────────────────────
\begin{document}

% ════════════════════════════════════════════════════════════
%  FRONT MATTER
% ════════════════════════════════════════════════════════════
\pagenumbering{roman}
\pagestyle{empty}

% ── Cover page ───────────────────────────────────────────────
\begin{titlepage}
\begin{center}
  \vspace*{0.5cm}

  % University logo placeholder — replace with:
  % \includegraphics[width=3.5cm]{images/asu_logo.png}
  \fbox{\parbox{3.5cm}{\centering\vspace{0.8cm}%
    \small[University Logo]\vspace{0.8cm}}}

  \vspace{0.6cm}
  {\large\textbf{Ain Shams University}}\\[2pt]
  {\large\textbf{Faculty of Computer \& Information Sciences}}\\[2pt]
  {\large\textbf{Computer Science Department}}

  \vspace{2.5cm}
  {\Huge\bfseries MyHR}\\[8pt]
  {\Large An Adaptive AI-Powered Hiring \& Interview Platform}

  \vspace{2cm}
  \textit{This documentation is submitted in partial fulfillment of the requirements\\
  for the Bachelor's degree in Computer and Information Sciences}

  \vspace{2cm}
  {\large\textbf{By:}}\\[10pt]
  \begin{tabular}{c}
    Pavlos Usama Youssef \quad [Computer Science]\\[4pt]
    Remon Osama Hosny Amin \quad [Computer Science]\\[4pt]
    Michael Ehad Karam \quad [Computer Science]\\[4pt]
    Mina Mossad Gerges \quad [Computer Science]\\[4pt]
    John Mahfouz Hanna \quad [Computer Science]
  \end{tabular}

  \vspace{1.8cm}
  {\large\textbf{Under Supervision of:}}\\[12pt]
  \textbf{Prof.\ Mohamad Mabrouk}\\
  {\small Lecturer in Computer Science Department,\\
  Faculty of Computer and Information Sciences, Ain Shams University.}\\[14pt]
  \textbf{TA.\ Manar Shaban}\\
  {\small Teaching Assistant in Computer Systems Department,\\
  Faculty of Computer and Information Sciences, Ain Shams University.}

  \vfill
  {\large June 2026}
\end{center}
\end{titlepage}

% ── Acknowledgements ─────────────────────────────────────────
\chapter*{Acknowledgements}
\addcontentsline{toc}{chapter}{Acknowledgements}

First and foremost, we would like to express our sincere gratitude to \textbf{Prof.\ Mohamad
Mabrouk} for his invaluable guidance, unwavering support, and insightful feedback throughout
the course of this graduation project. His mentorship has been instrumental in shaping both
the direction and the quality of our work.

We are equally grateful to \textbf{TA.\ Manar Shaban}, for her dedication, constructive input,
and the tremendous effort she invested in helping us navigate the challenges of this project.

We extend our deepest appreciation to each member of our \textbf{project team}. It was through
our collective dedication, seamless collaboration, and tireless effort that this vision was
successfully brought to life.

Furthermore, this endeavor would not have been possible without the robust academic foundation
and invaluable opportunities provided by the \textbf{Faculty of Computer and Information
Sciences} at \textbf{Ain Shams University}. We are profoundly proud to belong to an institution
whose environment consistently nurtures growth, innovation, and a genuine passion for our field.

Finally, to our esteemed professors, colleagues, friends, and families, thank you for your
endless patience and unwavering belief in us. Your support has been the cornerstone of our
success throughout this journey.

\clearpage

% ── Abstract ─────────────────────────────────────────────────
\chapter*{Abstract}
\addcontentsline{toc}{chapter}{Abstract}

Traditional recruitment and technical assessment processes often lack the adaptability required
to accurately evaluate candidates' technical competencies and soft skills in real time. Existing
interview platforms are typically limited by rigid assessment workflows, inefficient context
management, high computational costs, and limited scalability.

This project presents \textbf{MyHR}, an intelligent agentic platform that supports recruiters,
interviewers, and training evaluators through dynamic, context-aware interview simulations. The
platform leverages \textbf{Agentic Retrieval-Augmented Generation (Agentic RAG)} to maintain
conversational context, generate adaptive interview questions, and provide personalised
assessments based on candidate performance.

MyHR integrates a modular suite of artificial intelligence components, including automated CV
parsing, skill extraction and matching, semantic retrieval, cross-encoder reranking, and large
language model--based evaluation. A dynamic difficulty engine continuously adjusts question
complexity according to the candidate's responses, while multimodal analysis enhances assessment
by incorporating behavioural and non-verbal communication cues.

Experimental evaluation demonstrates that MyHR delivers accurate, scalable, and context-aware
candidate assessments while improving interview quality, reducing recruiter workload, and
providing meaningful insights for both enterprise hiring and technical training scenarios.

\vspace{1em}
\noindent\textbf{Keywords:} AI interview, retrieval-augmented generation, LLM agent, candidate
ranking, neural scoring, FastAPI, React, Firestore, adaptive difficulty.

\clearpage

% ── Table of Contents ────────────────────────────────────────
\pagestyle{fancy}
\tableofcontents
\clearpage

% ── List of Figures ──────────────────────────────────────────
\listoffigures
\addcontentsline{toc}{chapter}{List of Figures}
\clearpage

% ── List of Tables ───────────────────────────────────────────
\listoftables
\addcontentsline{toc}{chapter}{List of Tables}
\clearpage

% ── List of Abbreviations ────────────────────────────────────
\chapter*{List of Abbreviations}
\addcontentsline{toc}{chapter}{List of Abbreviations}

\begin{longtable}{@{}L{2.5cm}L{10cm}@{}}
\toprule
\textbf{Abbreviation} & \textbf{Meaning} \\
\midrule
\endhead
\bottomrule
\endfoot
API   & Application Programming Interface \\
ATS   & Applicant Tracking System \\
BM25  & Best Matching 25 (sparse ranking function) \\
CV    & Curriculum Vitae (résumé) \\
EAR   & Eye Aspect Ratio \\
HR    & Human Resources \\
IPR   & Iris Position Ratio \\
JD    & Job Description \\
LLM   & Large Language Model \\
MLP   & Multi-Layer Perceptron \\
NDCG  & Normalised Discounted Cumulative Gain \\
PII   & Personally Identifiable Information \\
PPO   & Proximal Policy Optimisation \\
RAG   & Retrieval-Augmented Generation \\
RBAC  & Role-Based Access Control \\
RL    & Reinforcement Learning \\
RRF   & Reciprocal Rank Fusion \\
SPA   & Single-Page Application \\
STT   & Speech-to-Text \\
TTS   & Text-to-Speech \\
VAD   & Voice Activity Detection \\
WS    & WebSocket \\
YOLO  & You Only Look Once \\
HTTP  & Hypertext Transfer Protocol \\
GUI   & Graphical User Interface \\
AI    & Artificial Intelligence \\
ML    & Machine Learning \\
NLP   & Natural Language Processing \\
\end{longtable}

\clearpage

% ════════════════════════════════════════════════════════════
%  MAIN MATTER
% ════════════════════════════════════════════════════════════
\pagenumbering{arabic}
\setcounter{page}{1}

% ════════════════════════════════════════════════════════════
%  CHAPTER 1 — INTRODUCTION
% ════════════════════════════════════════════════════════════
\chapterfrontpage{Chapter One}{Introduction}{%
  \item 1.1\quad Problem Definition
    \begin{itemize}
      \item 1.1.1\quad History
      \item 1.1.2\quad Applications
    \end{itemize}
  \item 1.2\quad Motivation
  \item 1.3\quad Objectives
  \item 1.4\quad Scope
  \item 1.5\quad Functional Requirements
  \item 1.6\quad Non-Functional Requirements
  \item 1.7\quad Project Timeline
  \item 1.8\quad Documentation Outline
}

\chapter{Introduction}
\label{ch:introduction}

\section{Problem Definition}

Hiring is one of the most resource-intensive processes a company undertakes, and two of its
earliest stages are also its most repetitive:

\begin{itemize}
  \item \textbf{Résumé screening.} For a single open role, recruiters may receive hundreds of
    CVs. Each must be read, compared against the job description, and judged for fit. The
    process is slow, subjective, and inconsistent between reviewers and across time of day.
  \item \textbf{First-round interviewing.} Even after shortlisting, conducting an initial
    screening interview for every promising candidate consumes scarce interviewer time.
    Scheduling alone introduces days of delay, and the quality of questions varies with the
    interviewer's preparation and fatigue.
\end{itemize}

These two stages share a common weakness: they depend on a human reading the same kinds of
documents and asking the same kinds of questions, over and over, with no guarantee of
consistency. The result is a hiring funnel that is \textbf{slow, expensive, hard to audit, and
vulnerable to unconscious bias}.

MyHR addresses this problem directly. It provides:

\begin{enumerate}
  \item An \textbf{enterprise portal} that ingests CVs in bulk, parses them, matches them
    against the job's required skills, and ranks candidates with neural models --- turning
    hundreds of documents into an ordered shortlist in seconds.
  \item An \textbf{AI interview engine} that conducts an adaptive, grounded interview with
    each invited candidate and returns a single, defensible score and a structured report,
    without occupying any human interviewer's time.
\end{enumerate}

\subsection{History}

The two enabling technologies behind MyHR matured only recently:

\begin{itemize}
  \item \textbf{Applicant Tracking Systems (ATS)} have existed for decades, but classical ATS
    rely on keyword matching. They reward candidates who echo the job description's wording
    rather than candidates who actually possess the underlying skills, and they cannot conduct
    an interview.
  \item \textbf{Large Language Models (LLMs)} became capable of fluent, context-aware question
    generation and answer evaluation only after the transformer architecture and
    instruction-tuned chat models. However, an LLM used naively will \textit{hallucinate} ---
    it may ask about experience the candidate never claimed.
  \item \textbf{Retrieval-Augmented Generation (RAG)} emerged as the standard technique to
    keep an LLM grounded in trustworthy source documents. By retrieving the most relevant
    passages of a CV and JD and feeding them to the model, RAG constrains the interview to
    facts that actually exist in the candidate's profile.
\end{itemize}

MyHR combines these threads: it replaces keyword ATS matching with neural skill matching, and
it replaces an ungrounded LLM with a RAG-grounded interview agent whose scores are
cross-checked by dedicated neural evaluators.

\subsection{Applications}

The platform serves two distinct audiences from one codebase:

\begin{itemize}
  \item \textbf{Enterprise hiring (primary).} A company requests access, is approved by a
    platform administrator, posts jobs, uploads candidate CVs, receives a neural ranking,
    invites the top candidates to an AI interview by email, and reviews completed interviews
    and aggregate analytics on a dashboard.
  \item \textbf{Candidate self-practice (secondary).} An individual candidate can sign up and
    run mock AI interviews against their own CV to prepare for real interviews.
\end{itemize}

Table~\ref{tab:comparison} positions MyHR against existing tools.

\begin{table}[H]
\centering
\caption{Comparison with Existing Solutions}
\label{tab:comparison}
\begin{tabularx}{\textwidth}{L{4.5cm}C{2.5cm}C{3.0cm}C{3.5cm}}
\toprule
\textbf{Capability} & \textbf{Keyword ATS} & \textbf{Generic LLM chatbot} & \textbf{MyHR} \\
\midrule
CV screening & Keyword match & None & Neural skill match + rubric \\
Conducts an interview & No & Yes (ungrounded) & Yes (CV/JD-grounded) \\
Answer scoring & None & Single opaque score & Transparent LLM + neural blend \\
Resists off-topic fluency & N/A & No & Yes (relevance gate) \\
Integrity / proctoring & No & No & Iris-gaze + face count \\
Multi-tenant enterprise workflow & Yes & No & Yes \\
Server-side, tamper-proof scoring & N/A & No & Yes \\
\bottomrule
\end{tabularx}
\end{table}

\section{Motivation}

The motivation for MyHR is to demonstrate that a \textbf{single, coherent system can automate
the top of the hiring funnel without sacrificing trustworthiness}. Three concerns drove the
design:

\begin{itemize}
  \item \textbf{Grounding over fluency.} An interviewer that sounds fluent but invents facts
    is worse than useless. MyHR's central engineering bet is that every question and every
    score must be grounded in retrieved evidence from the candidate's own CV and the JD.
  \item \textbf{Defensibility over convenience.} A candidate's score must be computed on the
    server from evidence the candidate cannot tamper with, and it must be explainable as a
    transparent blend of an LLM judgment and purpose-trained neural models.
  \item \textbf{Separation of concerns.} The enterprise workflow and the AI interview workflow
    are genuinely different problems, and the architecture keeps them in separate, well-defined
    layers.
\end{itemize}

\section{Objectives}

The project's objectives are grouped into five categories.

\textbf{Functional objectives} --- what the system must \textit{do}:

\begin{enumerate}
  \item Allow companies to onboard (request access $\to$ admin approval $\to$ invitation) and
    manage job postings.
  \item Ingest candidate CVs in bulk, parse them, and rank candidates against each job.
  \item Conduct an adaptive, voice-or-text AI interview with every invited candidate.
  \item Produce a per-candidate interview report and a combined hiring score, and present
    hiring analytics to HR users.
\end{enumerate}

\textbf{Technical objectives} --- how the system must be \textit{engineered}:

\begin{enumerate}
  \setcounter{enumi}{4}
  \item Implement a clean three-layer architecture (enterprise / interview-AI / training) with
    a FastAPI backend, a React SPA, and Firestore persistence.
  \item Expose a well-defined REST + WebSocket API and cover the core logic with automated
    tests.
  \item Keep the codebase maintainable and containerised for reproducible deployment.
\end{enumerate}

\textbf{AI objectives} --- the intelligence the system must exhibit:

\begin{enumerate}
  \setcounter{enumi}{7}
  \item Ground every interview question in retrieved evidence from the candidate's CV and JD
    (hybrid RAG) so the interviewer never hallucinates experience.
  \item Score answers through a transparent blend of an LLM judgment and purpose-trained
    neural models, guarded by a question-to-answer relevance gate.
  \item Train and evaluate the neural model layer rigorously on held-out test sets with
    standard ranking metrics, experiment tracking, and a human-rating validation study.
\end{enumerate}

\textbf{Security objectives}:

\begin{enumerate}
  \setcounter{enumi}{10}
  \item Authenticate all users via Firebase ID tokens and enforce RBAC separating candidates,
    HR users, and platform administrators.
  \item Compute candidate scores server-side only, and issue public interview links as
    cryptographically strong, expiring tokens.
  \item Redact PII before any document is placed in the vector store, and sanitise inputs
    against prompt-injection.
\end{enumerate}

\textbf{Performance objectives}:

\begin{enumerate}
  \setcounter{enumi}{13}
  \item Turn a batch of uploaded CVs into a ranked shortlist in seconds rather than hours.
  \item Serve hiring analytics from pre-computed aggregates and load neural models lazily.
\end{enumerate}

\section{Scope}

\textbf{In scope.} Enterprise account onboarding and approval; job posting; batch CV parsing,
skill extraction, and rubric scoring; neural candidate ranking; email-delivered, token-based
interview invitations; an adaptive, RAG-grounded, voice/text AI interview with silent
proctoring; server-side answer scoring and report synthesis; hiring analytics; candidate
self-practice; the full neural training and evaluation pipeline.

\textbf{Out of scope / partial.} Production-grade object storage is stubbed against local disk
and a MinIO/S3 configuration; observability is limited to a health endpoint, an in-process
metrics endpoint, and structured logging; the candidate ranker is trained on synthetically
generated comparative data pending real labelled outcomes; the human-rating validation study
was conducted with a single rater. These items are documented in Chapter~\ref{ch:conclusion}.

\section{Functional Requirements}

\begin{table}[H]
\centering
\caption{Functional Requirements}
\label{tab:fr}
\begin{tabularx}{\textwidth}{C{1.2cm}C{2.2cm}X}
\toprule
\textbf{ID} & \textbf{Actor} & \textbf{Requirement} \\
\midrule
FR-1  & HR (prospective) & Submit an enterprise access request using a corporate email address. \\
FR-2  & Admin            & Review, approve, or reject pending access requests. \\
FR-3  & HR               & Accept a company invitation and be granted the HR role. \\
FR-4  & HR               & Create, view, edit, close, and reopen job postings. \\
FR-5  & HR               & Upload candidate CVs in bulk (PDF/DOCX) for a job. \\
FR-6  & System           & Parse each CV, extract skills, match against the JD, and compute a rubric score. \\
FR-7  & HR               & View candidates ranked by match score and inspect per-candidate skill breakdown. \\
FR-8  & HR               & Invite a candidate to an AI interview by email link; blocked when job is closed. \\
FR-9  & Candidate        & Open a token-based interview link, grant camera/microphone access, take the interview. \\
FR-10 & System           & Generate adaptive, CV/JD-grounded questions; accept spoken or typed answers. \\
FR-11 & System           & Run silent proctoring (face count, gaze) and enforce anti-cheating timing. \\
FR-12 & System           & Score answers server-side and synthesise a per-candidate interview report. \\
FR-13 & HR               & View a candidate's interview report and combined hiring score; regenerate on demand. \\
FR-14 & HR               & View company-level hiring analytics and per-job pipeline statistics. \\
FR-15 & Candidate        & Self-register for candidate practice interviews. \\
\bottomrule
\end{tabularx}
\end{table}

\section{Non-Functional Requirements}

\begin{table}[H]
\centering
\caption{Non-Functional Requirements}
\label{tab:nfr}
\begin{tabularx}{\textwidth}{C{1.2cm}C{2.5cm}X}
\toprule
\textbf{ID} & \textbf{Category} & \textbf{Requirement} \\
\midrule
NFR-1  & Security        & All protected endpoints verify a Firebase ID token; scores are computed server-side only. \\
NFR-2  & Security        & Public interview links are cryptographically strong, single-purpose, and expire. \\
NFR-3  & Privacy         & PII is redacted before indexing; inputs are sanitised against prompt-injection. \\
NFR-4  & Multi-tenancy   & Every job-scoped request enforces company ownership (tenant isolation). \\
NFR-5  & Performance     & A batch of CVs is parsed and ranked in seconds; analytics served from pre-computed aggregates. \\
NFR-6  & Reliability     & Missing model checkpoints or external services degrade gracefully rather than crash. \\
NFR-7  & Usability       & Voice-native interview with clear status indicator, countdown, and accessible UI. \\
NFR-8  & Maintainability & Clean three-layer separation; core logic covered by an automated test suite. \\
NFR-9  & Portability     & Containerised with Docker; runs on CPU or CUDA GPU. \\
NFR-10 & Transparency    & Answer scores are an explainable weighted blend, not a single opaque number. \\
\bottomrule
\end{tabularx}
\end{table}

\section{Project Timeline}

The project was executed in five overlapping phases: foundations and research, core AI engine,
enterprise layer, training and evaluation, and hardening and documentation.

\begin{figure}[H]
\centering
\caption{Project Time Plan (Gantt)}
\label{fig:gantt}
\begin{tabularx}{\textwidth}{L{3.8cm}*{10}{C{0.85cm}}}
\toprule
\textbf{Phase} & \textbf{Oct} & \textbf{Nov} & \textbf{Dec} & \textbf{Jan} & \textbf{Feb}
              & \textbf{Mar} & \textbf{Apr} & \textbf{May} & \textbf{Jun} & \textbf{Jul} \\
\midrule
Foundations \& research  & \cellcolor{ainshams!60} & \cellcolor{ainshams!60} & & & & & & & & \\
Core AI engine           & & \cellcolor{ainshams!60} & \cellcolor{ainshams!60} & \cellcolor{ainshams!60} & & & & & & \\
Enterprise layer         & & & & \cellcolor{ainshams!60} & \cellcolor{ainshams!60} & \cellcolor{ainshams!60} & & & & \\
Training \& evaluation   & & & & & & \cellcolor{ainshams!60} & \cellcolor{ainshams!60} & & & \\
Hardening \& testing     & & & & & & & & \cellcolor{ainshams!60} & & \\
Documentation            & & \cellcolor{ainshams!30} & \cellcolor{ainshams!30} & \cellcolor{ainshams!30} & \cellcolor{ainshams!30} & \cellcolor{ainshams!30} & \cellcolor{ainshams!30} & \cellcolor{ainshams!30} & \cellcolor{ainshams!30} & \cellcolor{ainshams!30} \\
\bottomrule
\end{tabularx}
\end{figure}

\section{Documentation Outline}

The remainder of this document is organised as follows:

\begin{itemize}
  \item \textbf{Chapter~\ref{ch:background} (Background)} reviews the techniques and
    technologies the system relies on: retrieval-augmented generation, LLM agent
    orchestration, transformer embeddings and reranking, neural answer scoring, reinforcement
    learning for adaptive difficulty, emotion recognition and proctoring, and the platform
    stack.
  \item \textbf{Chapter~\ref{ch:design} (System Analysis \& Design)} presents the high-level
    architecture, the three-layer decomposition, the component and deployment diagrams, the
    database entity-relationship model, and the principal end-to-end workflows.
  \item \textbf{Chapter~\ref{ch:implementation} (System Implementation)} describes each
    component in detail --- the RAG pipeline, the interview agent, the neural models, the
    enterprise layer, proctoring and anti-cheating, the training layer, the database, the API
    surface, authentication and authorisation, and configuration.
  \item \textbf{Chapter~\ref{ch:testing} (System Testing)} explains the testing strategy, how
    to install, configure, and run the system, the automated test suite, and the end-to-end
    golden path.
  \item \textbf{Chapter~\ref{ch:results} (Results \& Discussion)} reports the measured
    outcomes --- model metrics, system performance, strengths, weaknesses, and lessons learned.
  \item \textbf{Chapter~\ref{ch:conclusion} (Conclusion and Future Work)} summarises the
    achievements, states the known limitations honestly, and proposes future improvements.
\end{itemize}

% ════════════════════════════════════════════════════════════
%  CHAPTER 2 — BACKGROUND
% ════════════════════════════════════════════════════════════
\chapterfrontpage{Chapter Two}{Background}{%
  \item 2.1\quad Retrieval-Augmented Generation
  \item 2.2\quad LLM Agent Orchestration
  \item 2.3\quad Transformer Embeddings \& Reranking
  \item 2.4\quad Neural Answer Scoring
  \item 2.5\quad Reinforcement Learning for Adaptive Difficulty
  \item 2.6\quad Emotion Recognition \& Proctoring
  \item 2.7\quad Platform Technologies
}

\chapter{Background}
\label{ch:background}

This chapter introduces the concepts, algorithms, and technologies that MyHR is built upon. It
is intended to give the reader the background necessary to follow the system design in
Chapter~\ref{ch:design} and the implementation in Chapter~\ref{ch:implementation}.

\section{Retrieval-Augmented Generation}

A \textbf{Large Language Model (LLM)} generates text by predicting the most likely continuation
of a prompt. While powerful, an LLM only ``knows'' what was in its training data and what is
placed in its prompt; asked about a specific candidate, it will confidently \textit{hallucinate}
details. To ground an LLM in trustworthy, up-to-date facts, \textbf{Retrieval-Augmented
Generation (RAG)} first \textit{retrieves} the most relevant passages from a source corpus and
then \textit{augments} the LLM's prompt with them, so the model reasons over real evidence
rather than guesses.

A retrieval system can be \textbf{sparse} or \textbf{dense}:

\begin{itemize}
  \item \textbf{Sparse retrieval} matches exact terms. \textbf{BM25} (Best Matching 25) is the
    classical sparse ranking function; it scores a document by the frequency of the query's
    terms within it, discounted by how common those terms are across the corpus. BM25 excels at
    exact keyword overlap (e.g.\ a specific framework name) but misses paraphrases.
  \item \textbf{Dense retrieval} matches \textit{meaning}. Each passage is encoded into a
    high-dimensional vector (an \textit{embedding}) by a transformer model, and relevance is
    measured by vector similarity (cosine distance). Dense retrieval captures semantic
    similarity even when no words overlap, but can miss rare exact terms.
\end{itemize}

Because the two approaches have complementary strengths, modern systems \textbf{fuse} them.
\textbf{Reciprocal Rank Fusion (RRF)} combines several ranked lists into one by summing, for
each document, a score of $\frac{1}{k + \text{rank}}$ across the lists, where $k$ is a
smoothing constant (commonly 60). Finally, a \textbf{cross-encoder reranker} re-scores the
fused shortlist by reading the query and each candidate passage \textit{together}, yielding the
most precise ordering at the cost of more computation.

MyHR uses exactly this stack: BM25 + dense (Pinecone) retrieval, RRF with $k = 60$, and a
cross-encoder reranker, to ground each interview question in the candidate's CV and the JD.

\section{LLM Agent Orchestration}

A single LLM call is stateless. A realistic interview, however, is a \textit{process}: rewrite
the context, retrieve evidence, check that the evidence is good enough, generate a question,
accept an answer, evaluate it, decide whether to continue, and adapt. Coordinating these steps
requires an \textbf{agent} --- a controller that maintains state and routes execution between
specialised steps.

\textbf{LangGraph} models such an agent as a \textbf{state graph}: a set of \textit{nodes}
(each a function that reads and writes a shared state object) connected by \textit{edges}. Some
edges are \textbf{conditional} --- the next node is chosen at runtime based on the current
state. A \textbf{checkpointer} persists the state between turns so a long-running,
multi-question interview can resume exactly where it left off.

MyHR's interview agent is a LangGraph state machine whose nodes are \texttt{rewrite},
\texttt{retrieve}, \texttt{grade}, \texttt{generate}, and \texttt{process\_answer}.

\section{Transformer Embeddings \& Reranking}

The quality of dense retrieval depends entirely on the embedding model. MyHR uses
\textbf{\texttt{all-mpnet-base-v2}}, a sentence-transformer that maps text to a
\textbf{768-dimensional} vector. It is used consistently at both indexing time and query time
--- a deliberate choice, because using one embedder to \textit{train} a model and a
\textit{different} one at inference silently degrades accuracy.

A \textbf{cross-encoder} differs from the bi-encoder used for retrieval. A bi-encoder embeds
the query and the document \textit{independently} and compares the two vectors --- fast, because
document vectors can be precomputed. A cross-encoder feeds the query and document
\textit{together} through the transformer and outputs a single relevance score --- slower, but
far more accurate. MyHR applies a cross-encoder only as a final reranking step over the small
fused candidate set.

\section{Neural Answer Scoring}

Relying on an LLM as the \textit{sole} judge of an answer has a weakness: the score is an
opaque opinion of one model. MyHR therefore complements the LLM with purpose-trained neural
networks:

\begin{itemize}
  \item A \textbf{Multi-Layer Perceptron (MLP)} is a feed-forward neural network of fully
    connected layers with non-linear activations. Given a fixed-length feature vector (here,
    embeddings of the question and answer), an MLP can be trained to regress a quality score.
  \item A \textbf{multi-head} network shares a common representation but ends in several
    independent output ``heads,'' each predicting a different target. MyHR's evaluator uses
    three heads to score an answer's \textbf{relevance}, \textbf{clarity}, and \textbf{depth}
    separately.
\end{itemize}

These models are evaluated with \textbf{rank-aware metrics} --- most importantly
\textbf{Spearman's rank correlation}, which measures whether the model orders answers the same
way the ground truth does, and \textbf{NDCG} (Normalised Discounted Cumulative Gain) for
ranking quality.

\section{Reinforcement Learning for Adaptive Difficulty}

A good interviewer adapts: if a candidate answers easily, the questions get harder; if they
struggle, the questions get easier. This is naturally framed as a \textbf{reinforcement learning
(RL)} problem, where an \textit{agent} observes a \textit{state} (the candidate's recent
performance), takes an \textit{action} (choose the next question's difficulty), and receives a
\textit{reward} (a more informative interview).

\textbf{Proximal Policy Optimisation (PPO)} is a widely used, stable policy-gradient RL
algorithm that improves the policy in small, clipped steps to avoid destructive updates.
MyHR's adaptive-difficulty module is trained in a simulated interview environment to map a
compact performance state to a difficulty action.

\section{Emotion Recognition \& Proctoring}

Two auxiliary signals enrich the interview:

\begin{itemize}
  \item \textbf{Emotion recognition} analyses the candidate's facial expression and/or vocal
    tone to estimate affective state during the interview, providing context for the report.
  \item \textbf{Proctoring} runs silently to detect integrity issues --- whether a face is
    present, whether the candidate is looking away, and whether multiple faces appear. The
    primary detector is \textbf{MediaPipe FaceMesh}, which returns 478 facial landmarks
    including the iris contours and therefore supports true eye-gaze estimation, with
    \textbf{OpenCV YuNet} as a lightweight fallback for face counting and head pose.
\end{itemize}

\section{Platform Technologies}

MyHR is assembled from established, production-grade components.

\begin{table}[H]
\centering
\caption{Platform Technology Stack}
\label{tab:techstack}
\begin{tabularx}{\textwidth}{L{2.2cm}L{4.5cm}X}
\toprule
\textbf{Layer} & \textbf{Technology} & \textbf{Role in MyHR} \\
\midrule
Frontend   & React 18 + Vite 5              & Single-page application for candidates and HR \\
Frontend   & Radix UI + Tailwind CSS        & Accessible component library and styling \\
Frontend   & React Router, TanStack Query   & Routing and server-state management \\
Backend    & FastAPI (Python)               & REST + WebSocket API server \\
Backend    & Uvicorn                        & ASGI server runtime \\
AI Agent   & LangGraph + LangChain         & Interview agent state machine \\
LLM        & Groq \texttt{llama-3.3-70b}   & Question generation and answer evaluation \\
RAG        & LlamaIndex + Pinecone          & Dense vector index and retrieval \\
RAG        & \texttt{rank\_bm25}            & Sparse BM25 retrieval \\
Embeddings & \texttt{all-mpnet-base-v2}     & 768-D text embeddings \\
Neural     & PyTorch                        & Eight custom models \\
Speech     & Deepgram SDK                   & Speech-to-text and text-to-speech \\
Proctoring & MediaPipe + OpenCV (YuNet)     & Silent face/iris-gaze detection \\
Privacy    & Microsoft Presidio             & PII detection and redaction \\
Auth       & Firebase Authentication        & Identity and ID-token verification \\
Database   & Google Cloud Firestore         & Multi-tenant document persistence \\
Email      & Resend API / Gmail SMTP        & Invitation and notification delivery \\
Training   & MLflow + TensorBoard           & Experiment tracking and metrics \\
Rate limit & SlowAPI                        & Per-route request throttling \\
\bottomrule
\end{tabularx}
\end{table}

% ════════════════════════════════════════════════════════════
%  CHAPTER 3 — SYSTEM ANALYSIS & DESIGN
% ════════════════════════════════════════════════════════════
\chapterfrontpage{Chapter Three}{System Analysis \& Design}{%
  \item 3.1\quad System Architecture
  \item 3.2\quad Three-Layer Architecture
  \item 3.3\quad Component Diagram
  \item 3.4\quad Enterprise Hiring Workflow
  \item 3.5\quad Candidate Interview Workflow
  \item 3.6\quad Authentication \& Authorisation Design
  \item 3.7\quad Database Design (Entity-Relationship Model)
  \item 3.8\quad Deployment Design
  \item 3.9\quad Interview Activity Diagram
}

\chapter{System Analysis \& Design}
\label{ch:design}

This chapter presents the analysis and design of MyHR: how the major components are organised,
how they communicate, how users and data flow through the system, and how the persistence,
security, and deployment concerns are structured.

\section{System Architecture}

MyHR follows a classic \textbf{client--server} architecture. A React single-page application
runs in the browser and communicates with a FastAPI backend over HTTP (REST) and, during a live
interview, over a WebSocket. The backend orchestrates four external/internal services: the LLM
(Groq), the vector database (Pinecone), the neural model layer (PyTorch, in-process), and the
persistence and identity services (Firestore and Firebase Authentication).

\begin{figure}[H]
\centering
\diagramplaceholder{%
  \textbf{System Architecture Diagram}\\
  Shows: Browser (React SPA) $\leftrightarrow$ FastAPI Backend (REST + WebSocket) $\leftrightarrow$
  External services (Groq LLM, Pinecone, Firestore, Firebase Auth, Deepgram, Email).\\
  The frontend never contacts Groq, Pinecone, or neural models directly.%
}
\caption{System Architecture Overview}
\label{fig:sysarch}
\end{figure}

The frontend never talks to Groq, Pinecone, or the neural models directly; all AI work is
mediated by the backend, which is also the only place candidate scores are computed.

\section{Three-Layer Architecture}

Conceptually, the system decomposes into three layers with distinct responsibilities,
lifecycles, and audiences.

\begin{figure}[H]
\centering
\diagramplaceholder{%
  \textbf{Three-Layer Architecture Diagram}\\
  Layer 1 (Enterprise): Accounts \& RBAC, Jobs \& CV upload, Neural ranking, Invitations \& analytics.\\
  Layer 2 (Interview-AI): RAG grounding, LangGraph agent, Answer scoring \& blend, Proctoring \& report.\\
  Layer 3 (Training): Data generation, Model training (8), Held-out evaluation, Human-rating study.\\
  Arrows: Enterprise invites candidates into Interview-AI; Training produces checkpoints for both.%
}
\caption{Three-Layer Architecture}
\label{fig:threelayer}
\end{figure}

\begin{itemize}
  \item The \textbf{Enterprise Layer} (\texttt{hr\_routes.py}, frontend HR pages) is
    multi-tenant, authenticated, and persisted in Firestore. It owns the hiring funnel.
  \item The \textbf{Interview-AI Layer} (\texttt{server.py}, \texttt{agent.py},
    \texttt{ingest.py}, \texttt{retriever.py}, \texttt{models/}) is stateful and real-time.
    It conducts interviews and scores answers.
  \item The \textbf{Training Layer} (\texttt{training/}) is offline. It generates data,
    trains the eight neural models, and evaluates them on held-out test sets. Its outputs are
    the model checkpoints consumed by the other two layers.
\end{itemize}

\section{Component Diagram}

The following diagram shows the principal backend modules and their dependencies.

\begin{figure}[H]
\centering
\diagramplaceholder{%
  \textbf{Component Diagram}\\
  Nodes: \texttt{server.py}, \texttt{hr\_routes.py}, \texttt{agent.py}, \texttt{ingest.py},
  \texttt{retriever.py}, \texttt{cv\_parser.py}, \texttt{models/registry.py},
  \texttt{firestore\_client.py}, \texttt{email\_service.py}, \texttt{prompts.py}.\\
  Dependencies: server $\to$ agent, hr, registry; hr $\to$ cv, registry, fc, email;
  agent $\to$ retr, registry, prompts; retr $\to$ ingest; ingest $\to$ registry.%
}
\caption{Component Diagram}
\label{fig:component}
\end{figure}

\section{Enterprise Hiring Workflow}

The enterprise workflow is the hiring funnel, from a company requesting access through to
reviewing completed interviews.

\begin{figure}[H]
\centering
\diagramplaceholder{%
  \textbf{Enterprise Hiring Workflow (Sequence Diagram)}\\
  Actors/participants: HR User, Platform Admin, React SPA, FastAPI (hr\_routes), Firestore, Email Service, Candidate.\\
  Phase 1 (Job Setup): request-access $\to$ notify admin $\to$ approve $\to$ invite HR $\to$ HR creates job $\to$ uploads CVs $\to$ ranking.\\
  Phase 2 (Interview): HR invites candidate $\to$ candidate completes interview $\to$ HR reviews report.%
}
\caption{Enterprise Hiring Workflow (Sequence)}
\label{fig:entseq}
\end{figure}

\section{Candidate Interview Workflow}

When a candidate opens their interview link, they enter a stateful, real-time session driven
by the LangGraph agent.

\begin{figure}[H]
\centering
\diagramplaceholder{%
  \textbf{Candidate Interview Workflow (Sequence Diagram)}\\
  Participants: Candidate, Interview Portal, FastAPI (server.py), LangGraph Agent, Hybrid RAG, Groq LLM, Neural Layer.\\
  Steps: validate token $\to$ start interview $\to$ index CV+JD $\to$ generate first question (grounded) $\to$
  loop [submit answer $\to$ evaluate $\to$ blend 65/20/15 $\to$ adapt difficulty $\to$ next question] $\to$
  synthesise report $\to$ thank-you screen.%
}
\caption{Candidate Interview Workflow (Sequence)}
\label{fig:candseq}
\end{figure}

Two design properties are visible here. First, the question is \textbf{grounded} before the
LLM is ever called. Second, the final score is \textbf{synthesised on the server}; the
candidate is shown only a thank-you screen, never their score.

\section{Authentication \& Authorisation Design}

Identity is delegated to \textbf{Firebase Authentication}. The browser authenticates the user
and receives a signed \textbf{ID token} (a JSON Web Token); every protected backend request
carries this token in an \texttt{Authorization: Bearer \ldots} header.

Authorisation is \textbf{role-based}. Three roles exist:

\begin{itemize}
  \item \textbf{Candidate} --- self-registerable; may run practice interviews.
  \item \textbf{HR (enterprise)} --- \textit{not} self-assignable; granted only by accepting a
    company invitation, which adds the user's UID to that company's \texttt{adminUIDs} array.
  \item \textbf{Platform administrator} --- a fixed allowlist email with platform privileges.
\end{itemize}

\begin{figure}[H]
\centering
\diagramplaceholder{%
  \textbf{Authentication \& Authorisation Flow (Sequence Diagram)}\\
  User $\to$ Firebase Auth (email/password or OAuth) $\to$ Signed ID token (JWT) $\to$
  FastAPI verifies token $\to$ Firestore resolves role (candidate $|$ hr $|$ admin) $\to$
  Role-gated response (or 401/403).%
}
\caption{Authentication \& Authorisation Flow}
\label{fig:auth}
\end{figure}

Making the HR role grantable only through an invitation --- rather than a self-service toggle
--- means a malicious user cannot escalate themselves into another company's data.

\section{Database Design (Entity-Relationship Model)}

MyHR persists all enterprise state in \textbf{Cloud Firestore}, a NoSQL document database.
The logical entities and their relationships are shown in Figure~\ref{fig:er}; the physical
collection layout and field-level detail appear in Section~\ref{sec:database}.

\begin{figure}[H]
\centering
\diagramplaceholder{%
  \textbf{Entity-Relationship Diagram}\\
  Entities: COMPANIES, JOBS, CANDIDATES, COMPANYSTATS, PENDINGREQUESTS, INVITATIONTOKENS, USERS.\\
  Relationships: Companies owns many Jobs; Jobs contains many Candidates; Companies has one CompanyStats;
  PendingRequests creates Companies; InvitationTokens scoped to Companies; Users adminUIDs in Companies.%
}
\caption{Database Entity-Relationship Diagram}
\label{fig:er}
\end{figure}

\section{Deployment Design}

In development the system runs as two processes --- the Vite dev server (port 8080) proxying
\texttt{/api} to the FastAPI backend (port 8000).

\begin{figure}[H]
\centering
\diagramplaceholder{%
  \textbf{Deployment Diagram}\\
  Browser $\to$ React SPA (Vite :8080) $\to$ FastAPI :8000 (Uvicorn ASGI).\\
  FastAPI $\to$ Groq LLM, Pinecone, Cloud Firestore, Deepgram STT/TTS, Resend/Gmail SMTP.\\
  Optional: Object storage (MinIO/S3), Redis cache.
  Firebase Auth accessed directly from both browser and backend.%
}
\caption{Deployment Diagram}
\label{fig:deploy}
\end{figure}

\section{Interview Activity Diagram}

Figure~\ref{fig:activity} shows the control flow of a single interview turn from the
candidate's perspective, including the proctoring and anti-cheating timing introduced in the
portal.

\begin{figure}[H]
\centering
\diagramplaceholder{%
  \textbf{Interview Turn (Activity Diagram)}\\
  Question ends $\to$ 3-2-1 countdown $\to$ Record (VAD on) $\to$ Spoke in 45 s?\\
  No $\to$ Submit empty (anti-cheat) $\to$ STT + evaluate.\\
  Yes $\to$ Silent after speech? $\to$ Still talking: back to Record.\\
  Silent $\to$ 3 s submit countdown $\to$ Resume / Submit $\to$ STT + evaluate.\\
  Evaluate $\to$ More questions? Yes: next turn. No: Report + thank-you.%
}
\caption{Interview Turn (Activity Diagram)}
\label{fig:activity}
\end{figure}

% ════════════════════════════════════════════════════════════
%  CHAPTER 4 — SYSTEM IMPLEMENTATION
% ════════════════════════════════════════════════════════════
\chapterfrontpage{Chapter Four}{System Implementation}{%
  \item 4.1\quad Hybrid RAG Pipeline
  \item 4.2\quad LangGraph Interview Agent
  \item 4.3\quad AI/ML Model Layer
  \item 4.4\quad Enterprise Layer
  \item 4.5\quad Training Layer
  \item 4.6\quad Database Design
  \item 4.7\quad API Reference
  \item 4.8\quad Authentication \& Authorisation
  \item 4.9\quad Configuration
  \item 4.10\quad Proctoring, Speech \& Anti-Cheating
  \item 4.11\quad Cross-Cutting Concerns
}

\chapter{System Implementation}
\label{ch:implementation}

This chapter describes the detailed implementation of each component using diagrams and prose
rather than source listings, with a short results-and-discussion note where relevant.

\section{Hybrid RAG Pipeline}

\textbf{Modules:} \texttt{ingest.py} (indexing) and \texttt{retriever.py} (retrieval).

When an interview starts, the candidate's CV and the JD are turned into a private, searchable
index scoped to the interview session. Before indexing, every document passes through two
safety filters: \textbf{PII redaction} (Microsoft Presidio) and a \textbf{prompt-injection
sanitiser}. The cleaned text is split into overlapping chunks (\texttt{chunk\_size = 512},
\texttt{chunk\_overlap = 50}), each CV chunk is prefixed with a small header carrying the
candidate name and role, and the chunks are written to two stores in parallel:
\textbf{Pinecone} (dense vectors, \texttt{all-mpnet-base-v2}, 768-D, namespaced by session)
and a \textbf{local BM25 store} (raw chunk text for sparse retrieval).

\begin{figure}[H]
\centering
\diagramplaceholder{%
  \textbf{Hybrid RAG Pipeline}\\
  CV + JD text $\to$ PII redaction + injection filter $\to$ Chunk 512 / overlap 50 + CV header $\to$
  Pinecone (dense 768-D) + BM25 (sparse) $\to$ RRF fusion ($k=60$) $\to$
  Cross-encoder rerank $\to$ Grounded passages.%
}
\caption{Hybrid RAG Pipeline (indexing and retrieval)}
\label{fig:rag}
\end{figure}

The hybrid design captures both exact-term matches (BM25) and semantic matches (dense
retrieval). The embedder is initialised lazily at startup so the $\sim$400\,MB model does not
block application import, and the same embedder is used at indexing and query time to avoid a
train/inference representation mismatch.

\section{LangGraph Interview Agent}

\textbf{Module:} \texttt{agent.py}. The interview is a LangGraph \textbf{state graph}
compiled with a checkpointer so each session's state survives between turns.

\begin{figure}[H]
\centering
\diagramplaceholder{%
  \textbf{LangGraph Agent State Graph}\\
  Start $\to$ rewrite $\to$ retrieve $\to$ grade (sufficient? $\to$ generate $\to$ emit question; insufficient? $\to$ rewrite).\\
  Answer received $\to$ process\_answer (continue $\to$ generate; end $\to$ finish).%
}
\caption{LangGraph Agent State Graph}
\label{fig:langgraph}
\end{figure}

\begin{itemize}
  \item \textbf{rewrite} --- reformulates the current topic into an effective retrieval query.
  \item \textbf{retrieve} --- runs the hybrid RAG pipeline (Section~\ref{sec:ragpipeline}) to
    gather grounded evidence.
  \item \textbf{grade} --- checks whether the retrieved context is good enough; a conditional
    edge loops back to re-query if not.
  \item \textbf{generate} --- prompts the Groq LLM (\texttt{llama-3.3-70b-versatile}) to
    produce the next question from the grounded context.
  \item \textbf{process\_answer} --- evaluates the candidate's answer, updates running
    performance, adapts difficulty, and a conditional edge decides whether to ask another
    question or end.
\end{itemize}

\subsection{Answer Scoring and the Blend}

When an answer is received, three independent quality signals are computed and combined
(Figure~\ref{fig:blend}). The \textbf{LLM judge} produces an overall score on a 0--100 scale;
the \textbf{MultiHeadEvaluator} (MOD-4) scores the answer's relevance, clarity, and depth; and
the \textbf{CandidateScoringMLP} (MOD-1) produces a deep-learned score from the question and
answer embeddings. Before blending, a \textbf{relevance gate} scales the neural contribution
by the cosine similarity between the question and answer embeddings.

\begin{figure}[H]
\centering
\diagramplaceholder{%
  \textbf{Answer Scoring and Blend}\\
  Answer $\to$ LLM judge (0--100) + Evaluator MOD-4 + Scorer MOD-1.\\
  MOD-4 and MOD-1 pass through Q$\leftrightarrow$A relevance gate.\\
  All three feed Weighted blend (65\,/\,20\,/\,15) $\to$ Answer score.%
}
\caption{Answer Scoring and Blend}
\label{fig:blend}
\end{figure}

The blend weights are fixed: when all three signals are available the score is
\textbf{65\% LLM + 20\% evaluator + 15\% MOD-1}; when MOD-1 is unavailable the system
degrades gracefully to \textbf{65\% LLM + 35\% evaluator}. Proctoring observations are
aggregated per answer and a flagged answer incurs a score penalty at this stage.

\section{AI/ML Model Layer}

\textbf{Modules:} \texttt{models/} with the lazy loader in \texttt{models/registry.py}. The
registry checks each checkpoint at startup (a health check), but loads each model into memory
only on first use.

\begin{table}[H]
\centering
\caption{Neural Model Layer}
\label{tab:models}
\begin{tabularx}{\textwidth}{C{0.6cm}L{3.2cm}L{2.5cm}L{2.3cm}L{2.0cm}X}
\toprule
\textbf{\#} & \textbf{Model} & \textbf{File} & \textbf{Input} & \textbf{Output} & \textbf{Role} \\
\midrule
MOD-1 & CandidateScoringMLP & \texttt{scoring\_model.py} & Q+A embeddings (1536-D) & Score 0--100 & Deep answer scoring in the blend \\
MOD-4 & MultiHeadEvaluator & \texttt{multi\_head\_evaluator.py} & Answer embed (768-D) & relevance, clarity, depth & Multi-dimensional answer evaluation \\
MOD-5 & NeuralCandidateRanker & \texttt{candidate\_ranker.py} & 7-D features & Ranking score & Order candidates within a job \\
MOD-6 & PerformancePredictor & \texttt{performance\_predictor.py} & Interview features & Market estimate & Auxiliary report signal \\
--- & SkillMatchSiamese & \texttt{skill\_matcher.py} & CV skills, JD skills & Match similarity & CV$\leftrightarrow$JD skill matching \\
--- & AdaptiveDifficulty & \texttt{difficulty\_engine.py} & Performance state & Difficulty action & RL-based question difficulty \\
--- & EmotionModel & \texttt{emotion\_model.py} & Face/tone features & Emotion estimate & Affective context for the report \\
--- & CrossEncoderScorer & \texttt{cross\_encoder\_scorer.py} & (query, passage) & Relevance & RAG reranking \\
\bottomrule
\end{tabularx}
\end{table}

\section{Enterprise Layer}

\textbf{Module:} \texttt{hr\_routes.py}, with \texttt{cv\_parser.py} and
\texttt{firestore\_client.py}. Key implementation details:

\begin{itemize}
  \item \textbf{Batch CV upload} parses each file, extracts skills, runs the neural
    \textbf{skill matcher} against the JD, computes a \textbf{rubric score}, and persists each
    candidate. Uploads are capped at 5\,MiB per file, and duplicate CVs are skipped by content
    hash and email.
  \item \textbf{Rubric scoring} scores a CV on four axes --- technical stack, architecture,
    experience, and preferred/extra signals --- and applies a \textbf{knock-out rule}: if a CV
    has zero skill overlap against a JD that lists at least five extractable skills, its score
    is capped at 40.
  \item \textbf{Atomic job statistics} are updated within a Firestore transaction so concurrent
    uploads do not lose updates.
  \item \textbf{Pre-computed analytics.} \texttt{GET /analytics} reads a pre-computed
    \texttt{CompanyStats} document with an in-memory time-to-live cache, avoiding an
    O(jobs $\times$ candidates) scan on every dashboard load.
  \item \textbf{Ownership checks.} Every job-scoped route verifies that the job belongs to the
    caller's company before returning data, enforcing tenant isolation.
\end{itemize}

\section{Training Layer}

\textbf{Module:} \texttt{training/}. The training layer is offline and produces the
checkpoints consumed by the other layers.

\begin{figure}[H]
\centering
\diagramplaceholder{%
  \textbf{Training Pipeline}\\
  Data generation (\texttt{generate\_*\_data.py}) $\to$ 70\,/\,15\,/\,15 split (train / val / test) $\to$
  Train (seeds, AdamW, cosine LR, early stop) $\to$ MLflow + TensorBoard logging $\to$
  Validate (select best epoch) $\to$ Held-out TEST evaluation $\to$
  Checkpoint $\to$ \texttt{models/checkpoints/} $\to$ Human-rating study (Spearman + Cohen's $\kappa$).%
}
\caption{Training Pipeline}
\label{fig:training}
\end{figure}

Each trainer shares the same hygiene: fixed random seeds, an AdamW optimiser with
cosine-annealing learning-rate schedule, early stopping with best-state restoration, dual
MLflow + TensorBoard logging, and rank-aware metrics (Spearman, NDCG, F1). Models are trained
on a \textbf{70/15/15} split and final numbers are reported on the \textbf{held-out test fold},
not the validation fold used for early stopping.

\begin{table}[H]
\centering
\caption{Representative Held-Out TEST Metrics}
\label{tab:trainmetrics}
\begin{tabularx}{\textwidth}{L{5cm}L{3.5cm}C{3cm}}
\toprule
\textbf{Model} & \textbf{Metric} & \textbf{Value} \\
\midrule
MultiHeadEvaluator (relevance) & Spearman's $\rho$ & $\approx 0.95$ \\
MultiHeadEvaluator (clarity)   & Spearman's $\rho$ & $\approx 0.95$ \\
MultiHeadEvaluator (depth)     & Spearman's $\rho$ & $\approx 0.95$ \\
CandidateScoringMLP            & MAE                & $\approx 0.067$ \\
CandidateScoringMLP            & Spearman's $\rho$ & $\approx 0.93$ \\
\bottomrule
\end{tabularx}
\end{table}

\section{Database Design}
\label{sec:database}

MyHR persists all enterprise state in \textbf{Cloud Firestore}. Table~\ref{tab:firestore}
documents the physical collection layout.

\begin{table}[H]
\centering
\caption{Firestore Collections}
\label{tab:firestore}
\begin{tabularx}{\textwidth}{L{3.5cm}L{4.0cm}X}
\toprule
\textbf{Collection} & \textbf{Purpose} & \textbf{Key fields} \\
\midrule
\texttt{Companies}                  & Tenant record        & \texttt{name}, \texttt{adminUIDs}, \texttt{domain}, \texttt{settings} \\
\texttt{Jobs}                       & Job postings         & \texttt{companyId}, \texttt{title}, \texttt{description}, \texttt{extractedSkills}, \texttt{stats} \\
\texttt{Jobs/\{id\}/Candidates}     & Candidates per job   & \texttt{name}, \texttt{email}, \texttt{matchScore}, \texttt{interviewScore}, \texttt{totalScore}, \texttt{interviewStatus} \\
\texttt{Users}                      & Portal role registry & \texttt{uid}, \texttt{role} (\texttt{candidate}) \\
\texttt{PendingRequests}            & Access requests      & \texttt{companyName}, \texttt{contactEmail}, \texttt{status} \\
\texttt{InvitationTokens}           & Tokens               & \texttt{type}, \texttt{companyId}, \texttt{jobId}, \texttt{candidateId}, \texttt{expiresAt} \\
\texttt{CompanyStats}               & Pre-computed analytics & \texttt{stats}, \texttt{monthly\_trends}, \texttt{updatedAt} \\
\bottomrule
\end{tabularx}
\end{table}

The total candidate score combines the CV match and the interview at
\textbf{40\% CV match + 60\% interview}.

\section{API Reference}

\begin{table}[H]
\centering
\caption{API Reference: Interview \& System Endpoints (\texttt{server.py})}
\label{tab:api1}
\begin{tabularx}{\textwidth}{C{1.4cm}L{5.5cm}C{1.8cm}X}
\toprule
\textbf{Method} & \textbf{Path} & \textbf{Auth} & \textbf{Purpose} \\
\midrule
POST & \texttt{/start\_interview} & Token & Start a practice interview session \\
POST & \texttt{/candidate-interview/\{token\}/start} & Public token & Start a token-based interview \\
WS   & \texttt{/ws/interview/\{session\_id\}} & Session & Live audio/video interview channel \\
GET  & \texttt{/end\_interview/\{session\_id\}} & Session & End an interview and finalise \\
POST & \texttt{/analyze\_frame} & Session & Proctoring: analyse a video frame \\
POST & \texttt{/submit\_answer} & Session & Submit and evaluate an answer \\
GET  & \texttt{/api/proctoring/\{session\_id\}} & Session & Retrieve proctoring results \\
POST & \texttt{/candidates/rank} & Internal & Rank candidates with neural ranker \\
GET  & \texttt{/health} & Public & Liveness/health check \\
GET  & \texttt{/metrics} & Public & In-process metrics \\
\bottomrule
\end{tabularx}
\end{table}

\begin{table}[H]
\centering
\caption{API Reference: Enterprise Endpoints (\texttt{hr\_routes.py})}
\label{tab:api2}
\begin{tabularx}{\textwidth}{C{1.4cm}L{5.5cm}C{1.8cm}X}
\toprule
\textbf{Method} & \textbf{Path} & \textbf{Auth} & \textbf{Purpose} \\
\midrule
POST   & \texttt{/request-access}                  & Public   & Submit enterprise access request \\
GET    & \texttt{/admin/pending-requests}           & Admin    & List pending requests \\
POST   & \texttt{/admin/accept-request/\{id\}}      & Admin    & Approve request, create company \\
POST   & \texttt{/admin/reject-request/\{id\}}      & Admin    & Reject a request \\
GET    & \texttt{/invite/\{token\}/validate}         & Public   & Validate a company-access invite \\
POST   & \texttt{/invite/\{token\}/accept}           & Firebase & Link user to company (grants HR) \\
POST   & \texttt{/jobs}                              & HR       & Create a job \\
GET    & \texttt{/jobs}                              & HR       & List company jobs \\
PATCH  & \texttt{/jobs/\{id\}}                       & HR       & Edit / close / reopen a job \\
POST   & \texttt{/jobs/\{id\}/upload-cvs}            & HR       & Batch upload + score CVs \\
GET    & \texttt{/jobs/\{id\}/candidates}            & HR       & List candidates \\
GET    & \texttt{/jobs/\{id\}/candidates/\{cid\}}    & HR       & Get a candidate with report \\
POST   & \texttt{/jobs/\{id\}/invite-interview/\{cid\}} & HR   & Email an interview invite \\
GET    & \texttt{/candidate-interview/\{token\}/validate} & Public & Validate interview token \\
POST   & \texttt{/candidate-interview/\{token\}/complete} & Public & Finalise interview \\
GET    & \texttt{/analytics}                         & HR       & Company hiring analytics \\
POST   & \texttt{/jobs/\{id\}/rank-candidates}       & HR       & Neural ranking of candidates \\
\bottomrule
\end{tabularx}
\end{table}

\section{Authentication \& Authorisation}

Identity is provided by \textbf{Firebase Authentication}; the backend verifies the Firebase
\textbf{ID token} on protected routes. Three roles exist:

\begin{itemize}
  \item \textbf{Candidate} --- self-registerable (\texttt{POST /user/role} accepts only
    \texttt{candidate}).
  \item \textbf{HR (enterprise)} --- \textit{not} self-assignable. Granted only by accepting a
    company invitation. An email already registered as a candidate cannot also become enterprise
    --- one email maps to one role.
  \item \textbf{Super-admin} --- a fixed allowlist email with platform privileges.
\end{itemize}

Public interview links use cryptographically strong tokens (\texttt{secrets.token\_urlsafe})
with an expiry and are validated server-side.

\section{Configuration}

\begin{table}[H]
\centering
\caption{Environment Variables}
\label{tab:envvars}
\begin{tabularx}{\textwidth}{L{5.5cm}X}
\toprule
\textbf{Variable} & \textbf{Purpose} \\
\midrule
\texttt{GROQ\_API\_KEY}                  & Groq LLM access (\texttt{llama-3.3-70b-versatile}) \\
\texttt{DEEPGRAM\_API\_KEY}              & Speech-to-text / text-to-speech \\
\texttt{PINECONE\_API\_KEY}              & Pinecone vector index \\
\texttt{FIREBASE\_SERVICE\_ACCOUNT\_PATH} & Firebase Admin SDK credentials \\
\texttt{VITE\_FIREBASE\_*}               & Frontend Firebase web config \\
\texttt{RESEND\_API\_KEY}                & Resend email transport \\
\texttt{SMTP\_USER}, \texttt{SMTP\_PASS} & Gmail SMTP transport (App Password) \\
\texttt{MYHR\_BASE\_URL}                 & Frontend base URL for email links \\
\texttt{AWS\_*}, \texttt{MINIO\_ENDPOINT} & Object storage (MinIO/S3) config \\
\texttt{REDIS\_URL}                      & Redis connection (caching) \\
\texttt{BYPASS\_EMAIL\_CHECK}            & Dev escape hatch for corporate-email gate \\
\bottomrule
\end{tabularx}
\end{table}

\section{Proctoring, Speech \& Anti-Cheating}

\subsection{Computer-Vision Proctoring}

Proctoring runs silently throughout the interview. For each captured video frame, the proctor
estimates three things: how many faces are present, whether the candidate is looking away, and
whether the eyes are closed. It uses a three-tier detection chain:

\begin{enumerate}
  \item \textbf{MediaPipe FaceMesh} (\texttt{refine\_landmarks=True}) provides 478 facial
    landmarks, including the iris contours (landmarks 468--477). From these the proctor
    computes an \textbf{Iris Position Ratio (IPR)} and an \textbf{Eye Aspect Ratio (EAR)}.
    The final gaze score is a weighted combination of iris position (0.55) and head pose (0.45).
  \item \textbf{YuNet} (\texttt{cv2.FaceDetectorYN}) provides fast, reliable face counting
    and a head-pose fallback when MediaPipe is unavailable.
  \item \textbf{Haar cascade} is the final fallback.
\end{enumerate}

\begin{figure}[H]
\centering
\diagramplaceholder{%
  \textbf{Proctoring Detection Pipeline}\\
  Video frame $\to$ Frame usable? No $\to$ Skip; Yes $\to$ FaceMesh 478 pts (YuNet fallback) $\to$
  gaze\_score (0.55$\times$iris + 0.45$\times$head) + Face count $\to$ Per-answer aggregation $\to$
  Violation thresholds? Yes $\to$ Suspicious; No $\to$ Clean.%
}
\caption{Proctoring Detection Pipeline}
\label{fig:proctor}
\end{figure}

An answer is flagged \textit{suspicious} when multiple faces appear, the face is absent for
more than 20\% of its frames, or the candidate looks away for more than 30\% of its frames.
A suspicious answer incurs a score penalty ($-5$, or $-8$ when multiple faces are detected).

\subsection{Speech (STT / TTS)}

Each question is synthesised to speech with \textbf{Deepgram TTS} and played to the candidate;
each spoken answer is captured in the browser, sent to the backend, and transcribed with
\textbf{Deepgram STT}. The transcript is what the LLM judge and the neural evaluators actually
score, so a candidate may answer by voice or by typing with identical downstream handling.

\subsection{Anti-Cheating in the Candidate Portal}

After a question's audio finishes, a mandatory \textbf{3-second countdown} is shown before the
microphone becomes active. The VAD's silence tolerance is set so that a natural $\sim$250\,ms
mid-sentence pause does \textbf{not} end the answer. When the candidate falls silent after
speaking, a \textbf{3-second submission countdown} appears. If the candidate stays completely
silent for 45 seconds, the portal submits an empty answer rather than waiting indefinitely.

\section{Cross-Cutting Concerns}

\begin{itemize}
  \item \textbf{Rate limiting:} SlowAPI provides per-route request throttling.
  \item \textbf{Graceful degradation:} A missing checkpoint, unavailable MOD-1, or absent
    MediaPipe wheel each fall back rather than crash.
  \item \textbf{Lazy model loading:} The registry verifies checkpoints at startup but loads
    each model into memory only on first use.
  \item \textbf{Prompt engineering:} Six prompt types are used --- First Question, Follow-up
    Question, Drill-down, Rubric Evaluation, Context Grading, and Report Synthesis. When CV
    context is weak or unavailable, the system generates questions based only on the JD,
    reducing hallucination risk.
\end{itemize}

% ════════════════════════════════════════════════════════════
%  CHAPTER 5 — SYSTEM TESTING
% ════════════════════════════════════════════════════════════
\chapterfrontpage{Chapter Five}{System Testing}{%
  \item 5.1\quad Testing Strategy
  \item 5.2\quad Installation
  \item 5.3\quad Running the System
  \item 5.4\quad Automated Test Suite
  \item 5.5\quad End-to-End Walkthrough
}

\chapter{System Testing}
\label{ch:testing}

This chapter explains the testing strategy, how to install, configure, and run MyHR, the
automated test suite, and the end-to-end golden path with screenshots.

\section{Testing Strategy}

MyHR is tested at several levels, chosen to give the most confidence per unit of effort.

\begin{table}[H]
\centering
\caption{Testing Strategy by Level}
\label{tab:teststrategy}
\begin{tabularx}{\textwidth}{L{2.8cm}L{4.5cm}X}
\toprule
\textbf{Level} & \textbf{What is tested} & \textbf{How} \\
\midrule
Unit        & CV parsing, skill extraction, rubric scoring, corporate-email gate, token generation & \texttt{pytest} (\texttt{test\_cv\_parser.py}, \texttt{test\_hr\_helpers.py}) \\
Integration & Cross-module enterprise helpers & \texttt{pytest} (\texttt{test\_integration.py}) with fake model registry fixture \\
API         & Endpoint contracts (status codes, auth gating, payload shape) & Manual + \texttt{pytest} \\
AI / model  & Model quality on held-out test folds (Spearman, NDCG, MAE) & \texttt{training/evaluate\_*.py} \\
RAG         & Retrieval grounding quality & \texttt{training/evaluate\_rag.py} \\
Frontend    & React component behaviour & Vitest + React Testing Library \\
Manual / UAT & End-to-end golden path & Human walkthrough (Section~\ref{sec:e2e}) \\
\bottomrule
\end{tabularx}
\end{table}

A deliberate design choice is that the \textbf{non-deterministic AI is evaluated
statistically} (on test sets, by correlation metrics) rather than with brittle exact-output
assertions, while the \textbf{deterministic logic} (parsing, scoring rules, auth gates) is
pinned with fast, exact unit tests.

\section{Installation}

\textbf{Prerequisites:} Python 3.11+, Node.js 18+ and npm, accounts/keys for Groq, Deepgram,
Pinecone, Firebase, and an email transport (Resend or Gmail App Password).

\begin{lstlisting}[language=bash, caption={Backend installation}]
# from the project root
python -m venv .venv
source .venv/bin/activate      # Linux/macOS
# .venv\Scripts\activate       # Windows
pip install -r requirements.txt
\end{lstlisting}

\begin{lstlisting}[language=bash, caption={Frontend installation}]
npm install
\end{lstlisting}

Create a \texttt{.env} file in the project root with the variables listed in
Table~\ref{tab:envvars}. At minimum set \texttt{GROQ\_API\_KEY}, \texttt{DEEPGRAM\_API\_KEY},
\texttt{PINECONE\_API\_KEY}, \texttt{FIREBASE\_SERVICE\_ACCOUNT\_PATH}, the
\texttt{VITE\_FIREBASE\_*} web config, an email transport, and \texttt{MYHR\_BASE\_URL}.

\section{Running the System}

\begin{lstlisting}[language=bash, caption={Starting backend and frontend}]
# Terminal 1 — backend (FastAPI on port 8000)
python server.py

# Terminal 2 — frontend (Vite dev server on port 8080)
npm run dev
\end{lstlisting}

On a healthy start, the backend logs model-registry health checks for all checkpoints, loads
the skill matcher, initialises the LlamaIndex embedder, readies the proctor, loads the emotion
model, and reports \textit{Application startup complete} with Uvicorn listening on
\texttt{http://0.0.0.0:8000}. The frontend serves the SPA at \texttt{http://localhost:8080}.

\section{Automated Test Suite}

\begin{lstlisting}[language=bash, caption={Running the backend test suite}]
python -m pytest tests/ -q
\end{lstlisting}

\begin{table}[H]
\centering
\caption{Automated Test Summary}
\label{tab:tests}
\begin{tabularx}{\textwidth}{L{5cm}C{1.5cm}X}
\toprule
\textbf{Test file} & \textbf{Tests} & \textbf{Coverage focus} \\
\midrule
\texttt{tests/test\_cv\_parser.py}    & 38 & Email/phone/name/skill extraction, skill matching, rubric scoring (incl.\ \texttt{framework\_cap} knock-out rule) \\
\texttt{tests/test\_hr\_helpers.py}   & 19 & Corporate-email gate, secure token generation, CV-text validation \\
\texttt{tests/test\_integration.py}   & 13 & Cross-module integration of enterprise helpers \\
\midrule
\textbf{Total}                        & \textbf{70} & All passing \\
\bottomrule
\end{tabularx}
\end{table}

The \texttt{tests/conftest.py} file provides shared fixtures (including a fake model registry)
so the suite runs without loading the heavy neural models.

\section{End-to-End Walkthrough}
\label{sec:e2e}

The following golden path exercises the full system.

\textbf{1. Request enterprise access.} A company submits the access form; a corporate email is
required.

\textbf{[Add Screenshot: Enterprise ``Request Access'' form]}

\textbf{2. Admin approval.} The platform admin reviews pending requests and approves one,
which creates the company, generates an invitation token, and emails the HR contact a link.

\textbf{[Add Screenshot: Admin pending-requests review screen]}

\textbf{3. Accept invitation.} The HR user opens the emailed link, creates their account, is
added to the company's \texttt{adminUIDs}, and lands on the HR dashboard as an enterprise user.

\textbf{[Add Screenshot: HR accepting the company invitation]}

\textbf{4. Create a job and upload CVs.} The HR user posts a job (skills are auto-extracted
from the JD) and uploads candidate CVs in bulk; each CV is parsed, skill-matched, and
rubric-scored.

\textbf{[Add Screenshot: Job management and bulk CV upload]}

\textbf{5. Review ranked candidates.} Candidates appear ranked by match score; the HR user can
open a candidate to see the match breakdown.

\textbf{[Add Screenshot: Ranked candidates list]}

\textbf{6. Invite a candidate to interview.} The HR user invites a candidate; an interview
link is emailed to the candidate.

\textbf{[Add Screenshot: Inviting a candidate to interview]}

\textbf{7. Candidate AI interview.} The candidate opens the link, grants camera/microphone
access, and completes an adaptive, grounded interview (voice or text). Proctoring runs
silently; the candidate never sees a score, only a thank-you screen on completion.

\textbf{[Add Screenshot: Candidate interview portal (camera + question)]}

\textbf{8. Hiring analytics.} Back on the dashboard, the HR user reviews aggregate analytics
--- candidate counts, average match and interview scores, monthly trends --- served from the
pre-computed \texttt{CompanyStats} document.

\textbf{[Add Screenshot: HR analytics dashboard]}

\textit{Expected outcome.} After completion, the candidate's record shows an
\texttt{interviewScore}, a synthesised \texttt{interviewReport}, and a \texttt{totalScore}
(40\% CV match + 60\% interview), and the analytics document refreshes in the background.
This confirms the full funnel operates end-to-end.

% ════════════════════════════════════════════════════════════
%  CHAPTER 6 — RESULTS & DISCUSSION
% ════════════════════════════════════════════════════════════
\chapterfrontpage{Chapter Six}{Results \& Discussion}{%
  \item 6.1\quad Model Results (Held-Out Test Sets)
  \item 6.2\quad RAG \& Grounding Results
  \item 6.3\quad System \& Backend Performance
  \item 6.4\quad Enterprise Funnel Results
  \item 6.5\quad Strengths
  \item 6.6\quad Weaknesses
  \item 6.7\quad Lessons Learned
}

\chapter{Results \& Discussion}
\label{ch:results}

This chapter reports the measured outcomes of the project and discusses what they mean. Model
metrics are reported on \textbf{held-out test folds} (the fold used neither for training nor
for early-stopping), so they reflect generalisation rather than memorisation.

\section{Model Results (Held-Out Test Sets)}

The neural layer was trained with a consistent \textbf{70\,/\,15\,/\,15} train/validation/test
split, fixed random seeds, an AdamW optimiser with cosine-annealing learning rate, and early
stopping with best-state restoration.

\begin{table}[H]
\centering
\caption{Model Test-Set Metrics}
\label{tab:modelmetrics}
\begin{tabularx}{\textwidth}{L{6cm}C{2cm}L{3.5cm}C{2cm}}
\toprule
\textbf{Model} & \textbf{ID} & \textbf{Primary metric} & \textbf{Value} \\
\midrule
MultiHeadEvaluator --- relevance head & MOD-4 & Spearman's $\rho$ & $\approx 0.95$ \\
MultiHeadEvaluator --- clarity head   & MOD-4 & Spearman's $\rho$ & $\approx 0.95$ \\
MultiHeadEvaluator --- depth head     & MOD-4 & Spearman's $\rho$ & $\approx 0.95$ \\
CandidateScoringMLP                   & MOD-1 & Mean Absolute Error & $\approx 0.067$ \\
CandidateScoringMLP                   & MOD-1 & Spearman's $\rho$ & $\approx 0.93$ \\
\bottomrule
\end{tabularx}
\end{table}

The high Spearman correlations indicate that both models \textbf{order} answers very close to
the reference labels --- which is the property that matters for ranking candidates. The scoring
MLP's low MAE ($\approx 0.067$ on a 0--1 scale) shows the absolute predictions are also
well-calibrated.

A \textbf{human-rating validation study} compared the system's blended scores against human
ratings, reporting a system-versus-human \textbf{Spearman's $\rho \approx 0.95$} and
\textbf{Cohen's $\kappa \approx 0.50$}. The strong rank correlation is encouraging; the
moderate $\kappa$ reflects that exact-category agreement is harder than rank agreement.

\begin{figure}[H]
\centering
\diagramplaceholder{%
  \textbf{Score Distribution \& Model Agreement}\\
  Left: Evaluator predicted-vs-reference scatter plot (Spearman $\rho \approx 0.95$).\\
  Right: Blend score-distribution histogram across all test answers.\\
  Export from TensorBoard / MLflow training run and replace this placeholder.%
}
\caption{Score Distribution \& Model Agreement}
\label{fig:scoredist}
\end{figure}

\section{RAG \& Grounding Results}

The retrieval layer fuses BM25 and dense (Pinecone, \texttt{all-mpnet-base-v2}) retrieval with
RRF ($k = 60$) and a cross-encoder reranker. The most important result is the
\textbf{grounding guarantee}: because the agent retrieves CV/JD evidence \textit{before} the
LLM is prompted, generated questions stay anchored to experience the candidate actually
claimed.

The \textbf{question-to-answer relevance gate} prevents a fluent but off-topic answer from
earning the full neural score. The gate uses the cosine similarity between the question and
answer embeddings, scaled so that a similarity of 0.5 or above counts as fully relevant.

RAG evaluation with \texttt{training/evaluate\_rag.py} showed high context recall across most
job roles tested, with an overall mean recall score of $\approx 0.98$.

\section{System \& Backend Performance}

The system meets its responsiveness objectives (NFR-5):

\begin{itemize}
  \item \textbf{CV screening} turns a batch of uploaded CVs into a ranked shortlist in seconds.
  \item \textbf{Analytics} are served from a pre-computed \texttt{CompanyStats} document with
    an in-memory time-to-live cache, avoiding an O(jobs $\times$ candidates) scan.
  \item \textbf{Model loading} is lazy: the registry verifies checkpoints at startup but loads
    each model into memory only on first use.
  \item \textbf{Graceful degradation} (NFR-6): a missing checkpoint, an unavailable MOD-1, or
    an absent MediaPipe wheel each fall back rather than crash.
\end{itemize}

Formal load and latency benchmarking of the interview WebSocket under concurrency was not
performed and is listed as future work (Chapter~\ref{ch:conclusion}).

\section{Enterprise Funnel Results}

End-to-end, the platform delivers the complete hiring funnel: a company can move from
\textit{requesting access} to a \textit{ranked, interviewed shortlist} without a human reading
a single CV or conducting a single first-round interview. Each candidate's record ends with a
total score, a detailed interview report, tone analysis, and proctoring observations.

\section{Strengths}

\begin{itemize}
  \item \textbf{Grounded interviewing:} Every question is backed by retrieved CV/JD evidence;
    the agent never asks about experience the candidate did not claim.
  \item \textbf{Transparent scoring:} The 65/20/15 blend and the relevance gate are explicit
    and auditable, not a single opaque number.
  \item \textbf{Complete funnel:} One platform covers the full hiring workflow from access
    request to scored interview report.
  \item \textbf{Strong model metrics:} Spearman $\approx 0.93$--$0.95$ on held-out test sets
    with consistent training hygiene.
  \item \textbf{Security:} Server-side scoring, token-based invitations, PII redaction, and
    RBAC make the platform defensible.
\end{itemize}

\section{Weaknesses}

\begin{itemize}
  \item \textbf{Validation breadth:} Answer-quality labels and the primary judge both originate
    from an LLM; the human-rating study currently uses a single rater.
  \item \textbf{Candidate ranker data:} Trained on synthetically generated comparative data;
    best used as a tiebreaker rather than an absolute measure.
  \item \textbf{Object storage:} Media and reports are stored against local disk and MinIO
    rather than a hardened cloud bucket.
  \item \textbf{Observability:} No full production monitoring, alerting, or distributed tracing.
\end{itemize}

\section{Lessons Learned}

\begin{itemize}
  \item Grounding is non-negotiable: an LLM interviewer without RAG produces generic,
    hallucinated questions that break user trust immediately.
  \item A relevance gate is essential: without it, a fluent off-topic answer earns a falsely
    high neural score.
  \item Lazy model loading is necessary: the $\sim$400\,MB embedder would make every test run
    slow if loaded eagerly.
  \item Tenant isolation must be designed in from the start; retrofitting ownership checks
    after the fact is fragile.
\end{itemize}

% ════════════════════════════════════════════════════════════
%  CHAPTER 7 — CONCLUSION AND FUTURE WORK
% ════════════════════════════════════════════════════════════
\chapterfrontpage{Chapter Seven}{Conclusion and Future Work}{%
  \item 7.1\quad Conclusion
  \item 7.2\quad Known Limitations
  \item 7.3\quad Future Work
}

\chapter{Conclusion and Future Work}
\label{ch:conclusion}

\section{Conclusion}

MyHR set out to automate the two most repetitive stages of hiring --- CV screening and
first-round interviewing --- without sacrificing trustworthiness, and it achieves that goal as
a working, end-to-end system.

The project delivers three cooperating subsystems: a \textbf{hybrid RAG layer} that grounds
every interview question in the candidate's own CV and the job description using BM25, dense
Pinecone retrieval, Reciprocal Rank Fusion, and cross-encoder reranking; a \textbf{LangGraph
LLM agent} that conducts an adaptive, voice-or-text interview and evaluates answers through a
transparent blend of an LLM judgment and purpose-trained neural models (65\% LLM / 20\%
evaluator / 15\% deep scorer, with a relevance gate); and a \textbf{neural layer of eight
PyTorch models} trained with professional hygiene --- fixed seeds, early stopping, experiment
tracking, and reporting on held-out test sets --- that reach Spearman correlations around
0.93--0.95 against their reference labels.

Around this AI core, the system provides a complete \textbf{multi-tenant enterprise platform}:
Firebase-backed authentication, role-based access control, batch CV upload with neural skill
matching and rubric scoring, email-delivered token-based interview invitations,
server-side-only scoring that candidates cannot tamper with, and a pre-computed analytics
dashboard. The frontend is a polished React single-page application, the backend a
well-structured FastAPI service, and the whole is covered by an automated test suite (70
passing backend tests) and this documentation.

\textbf{Principal achievements:}

\begin{itemize}
  \item A grounded, adaptive AI interview engine with transparent, defensible scoring.
  \item Automated CV screening with neural skill matching and a rubric with a knock-out rule.
  \item A secure, multi-tenant enterprise portal with RBAC and token-based invitations.
  \item Eight neural models trained and evaluated on held-out test sets with experiment
    tracking.
  \item A clean separation of enterprise, interview-AI, and training layers, with automated
    tests.
\end{itemize}

\section{Known Limitations}

\begin{itemize}
  \item \textbf{Validation breadth.} Answer-quality labels and the primary judge both originate
    from an LLM; the human-rating validation study currently uses a single rater.
  \item \textbf{Candidate ranker data.} Trained on synthetically generated comparative data;
    its scores are best used as a tiebreaker.
  \item \textbf{Object storage.} Media and reports stored against local disk and MinIO rather
    than a hardened, durable cloud bucket.
  \item \textbf{Observability.} The backend exposes a health endpoint and structured logging,
    but does not yet include full production monitoring, alerting, or distributed tracing.
\end{itemize}

\section{Future Work}

\begin{itemize}
  \item \textbf{Multi-rater validation.} Extend the human-rating study to several independent
    raters and report inter-rater reliability alongside the system-vs-human correlation.
  \item \textbf{Real comparative labels for ranking.} Retrain the candidate ranker on real
    hiring outcomes to make its scores absolute rather than relative.
  \item \textbf{Durable object storage.} Replace the local-disk path with a hardened cloud
    object store for CVs, audio, and reports.
  \item \textbf{Production observability.} Add metrics dashboards, alerting, request tracing,
    and model drift monitoring.
  \item \textbf{Model registry and versioning.} Promote the file-based checkpoints to a
    versioned model registry with rollback.
  \item \textbf{Scaling and load testing.} Benchmark the interview WebSocket and batch upload
    paths under load, and introduce horizontal scaling.
  \item \textbf{Richer analytics and fairness reporting.} Surface the existing fairness-audit
    harness in the dashboard and expand analytics with funnel and time-to-hire metrics.
  \item \textbf{Personalised Career Coaching.} Add a career coaching module that recommends
    learning paths, projects, and courses based on the candidate's interview performance and
    skill gaps.
  \item \textbf{Collaborative Recruitment Workspace.} Add shared workspaces where multiple
    recruiters can review candidates, add comments, compare reports, and collaborate on final
    hiring decisions.
\end{itemize}

% ════════════════════════════════════════════════════════════
%  BACK MATTER
% ════════════════════════════════════════════════════════════

% ── Tools ────────────────────────────────────────────────────
\chapter*{Tools Used}
\addcontentsline{toc}{chapter}{Tools Used}

\begin{longtable}{L{2.8cm}L{4.5cm}X}
\caption{Tools Used} \label{tab:tools} \\
\toprule
\textbf{Category} & \textbf{Tool} & \textbf{Use in MyHR} \\
\midrule
\endfirsthead
\multicolumn{3}{l}{\small\textit{(continued from previous page)}} \\
\toprule
\textbf{Category} & \textbf{Tool} & \textbf{Use in MyHR} \\
\midrule
\endhead
\midrule
\multicolumn{3}{r}{\small\textit{(continued on next page)}} \\
\endfoot
\bottomrule
\endlastfoot
Language (backend)   & Python 3.14                             & Backend, AI engine, training \\
Language (frontend)  & JavaScript (ES2022)                     & React SPA \\
Web framework        & FastAPI + Uvicorn                        & REST + WebSocket API \\
Frontend build       & Vite 5                                  & Dev server and production bundling \\
UI                   & React 18, Radix UI, Tailwind CSS        & Component-based interface \\
Charts / UX          & Recharts, Framer Motion                 & Analytics charts and animations \\
AI agent             & LangGraph, LangChain                    & Interview state machine \\
LLM                  & Groq \texttt{llama-3.3-70b-versatile}  & Question generation, answer judging \\
Vector DB            & Pinecone (via LlamaIndex)               & Dense retrieval \\
Sparse retrieval     & \texttt{rank\_bm25}                     & BM25 retrieval \\
Embeddings           & \texttt{sentence-transformers}          & 768-D text embeddings \\
Deep learning        & PyTorch                                 & Eight neural models \\
Speech               & Deepgram SDK                            & STT / TTS \\
Computer vision      & MediaPipe FaceMesh + OpenCV (YuNet)     & Silent iris-gaze proctoring \\
Privacy              & Microsoft Presidio                      & PII redaction \\
Auth                 & Firebase Authentication                 & Identity, ID tokens \\
Database             & Google Cloud Firestore                  & Persistence \\
Email                & Resend API / Gmail SMTP                 & Invitations, notifications \\
Experiment tracking  & MLflow, TensorBoard                     & Training metrics \\
Rate limiting        & SlowAPI                                 & Per-route throttling \\
Testing              & pytest, Vitest, React Testing Library   & Automated tests \\
Version control      & Git / GitHub                            & Source management \\
IDE                  & Visual Studio Code                      & Development \\
Containerisation     & Docker + docker-compose                 & Reproducible deployment \\
\end{longtable}

\textbf{Hardware.} Development and training used a CUDA-capable GPU workstation; the system
also runs on CPU.

% ── References ───────────────────────────────────────────────
\chapter*{References}
\addcontentsline{toc}{chapter}{References}

\begin{enumerate}[leftmargin=2em]
  \item A. Vaswani, N. Shazeer, N. Parmar, et al., ``Attention Is All You Need,''
    \textit{Advances in Neural Information Processing Systems (NeurIPS)}, 2017.
  \item S. Robertson and H. Zaragoza, ``The Probabilistic Relevance Framework: BM25 and
    Beyond,'' \textit{Foundations and Trends in Information Retrieval}, 2009.
  \item N. Reimers and I. Gurevych, ``Sentence-BERT: Sentence Embeddings using Siamese
    BERT-Networks,'' \textit{EMNLP}, 2019.
  \item G. V. Cormack, C. L. A. Clarke, and S. Büttcher, ``Reciprocal Rank Fusion Outperforms
    Condorcet and Individual Rank Learning Methods,'' \textit{SIGIR}, 2009.
  \item P. Lewis, E. Perez, A. Piktus, et al., ``Retrieval-Augmented Generation for
    Knowledge-Intensive NLP Tasks,'' \textit{NeurIPS}, 2020.
  \item J. Schulman, F. Wolski, P. Dhariwal, A. Radford, and O. Klimov, ``Proximal Policy
    Optimization Algorithms,'' \textit{arXiv:1707.06347}, 2017.
  \item K. Song, X. Tan, T. Qin, J. Lu, and T.-Y. Liu, ``MPNet: Masked and Permuted
    Pre-training for Language Understanding,'' \textit{NeurIPS}, 2020.
  \item LangGraph Documentation.
    \url{https://langchain-ai.github.io/langgraph/}, last retrieved 27/06/2026.
  \item FastAPI Documentation.
    \url{https://fastapi.tiangolo.com/}, last retrieved 27/06/2026.
  \item React Documentation.
    \url{https://react.dev/}, last retrieved 27/06/2026.
  \item PyTorch Documentation.
    \url{https://pytorch.org/docs/stable/index.html}, last retrieved 27/06/2026.
  \item Pinecone Documentation.
    \url{https://docs.pinecone.io/}, last retrieved 27/06/2026.
  \item Groq Documentation.
    \url{https://console.groq.com/docs}, last retrieved 27/06/2026.
  \item Firebase Authentication \& Cloud Firestore Documentation.
    \url{https://firebase.google.com/docs}, last retrieved 27/06/2026.
  \item Deepgram Speech-to-Text \& Text-to-Speech Documentation.
    \url{https://developers.deepgram.com/}, last retrieved 27/06/2026.
  \item MediaPipe Face Mesh Documentation.
    \url{https://developers.google.com/mediapipe}, last retrieved 27/06/2026.
  \item Microsoft Presidio Documentation.
    \url{https://microsoft.github.io/presidio/}, last retrieved 27/06/2026.
\end{enumerate}

% ── Appendices ───────────────────────────────────────────────
\appendix

\chapter{Environment Variable Template}
\label{app:env}

\begin{lstlisting}[language=bash, caption={.env template}]
# LLM & AI services
GROQ_API_KEY=
DEEPGRAM_API_KEY=
PINECONE_API_KEY=
OPENAI_API_KEY=

# Firebase (Admin SDK + web config)
FIREBASE_SERVICE_ACCOUNT_PATH=
VITE_FIREBASE_API_KEY=
VITE_FIREBASE_AUTH_DOMAIN=
VITE_FIREBASE_PROJECT_ID=
VITE_FIREBASE_STORAGE_BUCKET=
VITE_FIREBASE_MESSAGING_SENDER_ID=
VITE_FIREBASE_APP_ID=

# Email transport (use one)
RESEND_API_KEY=
SMTP_USER=
SMTP_PASS=

# URLs & infrastructure
MYHR_BASE_URL=http://localhost:8080
AWS_ACCESS_KEY_ID=
AWS_SECRET_ACCESS_KEY=
AWS_REGION=
S3_BUCKET_NAME=
MINIO_ENDPOINT=
REDIS_URL=
DATABASE_URL=

# Dev escape hatch (omit in production)
BYPASS_EMAIL_CHECK=
\end{lstlisting}

\chapter{Repository Structure}
\label{app:repo}

\begin{lstlisting}[caption={Selected repository structure}]
MyHR/
|-- server.py              FastAPI app, interview WebSocket, system endpoints
|-- hr_routes.py           Enterprise router (jobs, CVs, invites, analytics)
|-- agent.py               LangGraph interview agent + scoring blend
|-- ingest.py              RAG indexing + lazy embedder
|-- retriever.py           BM25 + dense + RRF + cross-encoder rerank
|-- cv_parser.py           CV parsing, skill extraction, rubric scoring
|-- email_service.py       Resend / Gmail SMTP transport + templates
|-- firestore_client.py    Firestore helpers + role sync
|-- prompts.py             LLM prompt templates
|-- models/
|   |-- registry.py
|   |-- multi_head_evaluator.py
|   |-- scoring_model.py
|   |-- candidate_ranker.py
|   |-- skill_matcher.py
|   |-- difficulty_engine.py
|   |-- emotion_model.py
|   |-- performance_predictor.py
|   |-- cross_encoder_scorer.py
|   |-- proctor.py
|   `-- checkpoints/       Trained model weights
|-- training/              Data generation, trainers, evaluation, human study
|-- tests/                 pytest suite (70 tests)
|-- src/                   React SPA (pages, components, contexts, hooks)
`-- docs/                  This documentation set
\end{lstlisting}

\end{document}
